\section{Introduction}\label{sec:1}

For many job seekers, finding work still feels like shouting into a crowd.
They rewrite the same resume for dozens of roles, wait weeks without feedback, and rarely learn why one posting surfaced in their results while another vanished after a single click.
When the process does move forward, it often moves on the employer's timetable---not the candidate's---leaving people to guess what mattered most: their skills, their location, their salary expectations, or something they never saw in the job ad.

Recruitment systems were built to help organizations sort large applicant pools, but the experience on both sides remains uneven.
Candidates see opaque lists and generic rejections.
Recruiters and hiring managers see dashboards that prioritize speed over clarity: filters that drop promising profiles, rankings that change without explanation, and little visibility into how a job requirement maps to a person's actual background.
Trust erodes when neither party can answer a simple question---\emph{why this match, and why now?}

The cost of getting filtering wrong is measured in time and money, not just frustration.
Every unsuitable interview consumes calendar space, interviewer attention, and candidate hope.
Every missed fit---someone qualified who never appeared on a shortlist---can mean a longer vacancy, a rushed hire, or lost talent to a competitor.
At scale, small errors in matching compound into weeks of delay and substantial recruiting spend, even before anyone discusses compensation or start dates.

A broader vision of computing may help reframe what hiring tools could become.
In his writing on education and artificial intelligence, Sal Khan describes a future in which personal AI assistants walk alongside individuals---helping them understand options, prepare materials, and navigate complex decisions rather than replacing human judgment.
That idea transfers naturally to the job market: a job seeker deserves an assistant that represents \emph{their} goals, history, and constraints, not merely a form field on an employer's website.

We take that vision seriously by giving each side of the market its own software agent.
On the candidate side, a \textbf{client-side agent} accepts resumes and profile updates, helps organize skills and experience into a structured record, and keeps that record current as the person edits their materials.
On the employer side, an \textbf{employer-side agent} accepts job descriptions and hiring requirements, structures them into a comparable record, and maintains them as roles evolve.
Neither agent exists in isolation: a candidate's story and an employer's needs only become meaningful when they can be compared fairly, explained clearly, and updated when either party changes direction.

Collaboration among agents is therefore the design center, not an afterthought.
If candidate information and job information live in separate silos with no shared language, both portals revert to blind search.
If a single monolithic program owns every edit and every score, users cannot tell who changed what or why a ranking shifted.
JobMatch instead assigns clear ownership---candidates control candidate state, employers control job state---and introduces a dedicated \textbf{matchmaking agent} that reads both sides, scores possible pairings, and returns ranked recommendations with reasons attached.
When a resume or job description changes, the owning agent updates its record and notifies the matchmaking layer so suggestions stay fresh.

Together, these pieces form \textbf{JobMatch}, a multi-agentic job-matching system implemented as a research prototype with separate web experiences for candidates, employers, and administrators.
The system parses resumes and job descriptions into structured states, builds comparable representations of people and roles, and routes matching work through the matchmaking agent.
Users remain in charge: they review parsed fields, inspect why a pairing scored as it did, and explicitly choose to save, apply, or shortlist---the software suggests; it does not hire on anyone's behalf.
We evaluate the matching core offline on a labeled demo corpus (30 resumes, 15 jobs, 47 graded relevance pairs) using standard quality metrics and reproducible benchmark comparisons; we report numbers only from committed artifacts and do not claim production hiring effectiveness.

This paper makes the following contributions:
\begin{itemize}[noitemsep]
  \item A \textbf{candidate/client-side agent} that processes resumes and CVs and maintains candidate state.
  \item An \textbf{employer-side agent} that processes job descriptions and maintains job state.
  \item A \textbf{matchmaking engine} that uses semantic matching and ranking to produce ordered recommendations.
  \item A \textbf{multi-agent communication flow} in which agents share relevant state and data when profiles change.
  \item A \textbf{complete UI and application layer} demonstrating end-to-end candidate, employer, and admin workflows.
  \item An \textbf{evaluation} using quality metrics and benchmark comparisons on a fixed labeled corpus.
\end{itemize}

The remainder of the paper is organized as follows.
Section~\ref{sec:2} reviews recruitment systems, AI agents in hiring, and semantic matching.
Section~\ref{sec:3} presents the architecture of the multi-agent system.
Section~\ref{sec:4} describes implementation details.
Section~\ref{sec:5} defines quality metrics and the evaluation protocol.
Section~\ref{sec:6} reports results and discussion.
Section~\ref{sec:7} concludes and outlines future scope.
